% \section{Eksisterende videoløsninger}
% \label{sec:eksisterende}

% HDO leverer i dag to separate videoløsninger til den akuttmedisinske kjede.
% Begge bruker \gls{webrtc} for å strømme video fra innringers kamera til
% operatøren, og begge følger samme arbeidsflyt: operatøren sender en
% SMS-lenke til innringeren, som godtar kameradeling og starter
% videostrømmen. Videostrømmen er kryptert og går via en sentralisert
% videomotor, mens lyd overføres via den ordinære telefonsamtalen
% \cite{hdo_oppgave}.

% Til tross for lik funksjonalitet benytter de to løsningene forskjellig
% infrastruktur og videomotorer: AntMedia Server og NHN Video. Hver motor
% har sitt eget API og sine egne integrasjonsmetoder, noe som innebærer at
% ny funksjonalitet må implementeres separat for begge, og at driftspersonell
% må forholde seg til to ulike systemer \cite{hdo_oppgave}.

% VideoAPI-prosjektet ble initiert for å adressere denne dupliseringen ved
% å tilby én mellomvare som abstraherer forskjellene mellom motorene, slik
% at fremtidige endringer kan gjøres ett sted.

\section{Eksisterende videoløsninger}
\label{sec:eksisterende}
HDO leverer i dag to separate videoløsninger til den akuttmedisinske kjede. Begge følger arbeidsflyten beskrevet i delkapittel~\ref{sec:hdo} og bruker \gls{webrtc} for å strømme kryptert video fra innringers kamera til operatøren \cite{hdo_oppgave}.

Til tross for lik funksjonalitet er de to løsningene basert på forskjellige videomotorer med ulik underliggende teknologi: AntMedia Server og NHN Video. De to motorene beskrives nærmere i avsnitt~\ref{sec:antmedia} og~\ref{sec:nhn}.